\documentclass{article}

\usepackage{french}
 
\pagestyle{headings}
\markright{28/6/1995 \hfill Conservatoire National des Arts et
M\'etiers -- Paris \hfill\ } 

\baselineskip 6pt
\parskip 6pt
\onecolumn
\textwidth 16cm
\textheight 24cm
\hoffset=-2cm
\voffset=-2.5cm
\setlength{\parindent}{0.0cm}
\unitlength 1cm
\input{psfig}
\newtheorem{defin}{D\'efinition}[section]
\newtheorem{exampl}{Exemple}[section]
\def\picwidth{16cm}

\begin{document}
\psfig{figure=dept.epsi} 
\begin{center}
{\Large\bf Solution du  Partiel Bases de Donn\'ees}
\end{center}

%%%%%%%%%%%%%%%%%%%%%%%%%%%%%%%%%%%%%%%%%%%%%%%%%%%%%%%%%%%%

\vspace*{2cm}

\smallskip
{\bf Cr\'eation du sch\'ema}

Il y a 5 tables\,: une pour chaque entit\'e, et une pour
      l'association n-m ``fr\'equente''. Voici le sch\'ema, avec
       les contraintes.

\begin{verbatim}
CREATE TABLE Hotel (nomHotel VARCHAR2 (30),
	            nomGerant      VARCHAR2(30) NOT NULL,
                    etoiles      NUMBER (10,0),
	             PRIMARY KEY (nomHotel), 
                      UNIQUE (nomGerant),
                    );

CREATE TABLE Chambre (nomHotel   VARCHAR2 (30),
                       No      NUMBER(5),
                       capacite         NUMBER (2) NOT NULL,
                       PRIMARY KEY (nomHotel, No),
                       FOREIGN KEY (nomHotel) REFERENCES Hotel
                         ON DELETE CASCADE
                       );
\end{verbatim}
Hotel/Chambre est une association de composition\,: il
est donc normal d'identifier une chambre par
le couple {\tt (nomHotel, no)}. Consid\'erer que le
{\tt No} est l'identifiant de la chambre est acceptable,
mais complique les requ\^etes ensuite.

\begin{verbatim} 
CREATE TABLE Skieur (id NUMBER (10),
                      nomSkieur VARCHAR2(30) NOT NULL,
                       nomHotel VARCHAR2(30),
                       noChambre NUMBER (5),
                       PRIMARY KEY (id),
                       FOREIGN KEY (nomHotel, noChambre) REFERENCES Chambre
                       ON DELETE SET NULL
                       CHECK 0 <= (SELECT MAX(capacite) - COUNT(*)
                                   FROM   Chambre c, Skieur s
                                   WHERE  c.no = s.noChambre
                                   GROUP BY noChambre)               
                   );  
\end{verbatim}
La commande {\tt CHECK} ci-dessus \'etait franchement difficile...

\begin{verbatim}
CREATE TABLE Piste  (id NUMBER (10),
                       nomPiste VARCHAR2 (30),
                       couleur VARCHAR2(30) NOT NULL,
                        PRIMARY KEY (id),
                        CHECK (couleur IN ('Rouge', 'Bleue', 'Noire',
                                             'Verte')));

CREATE TABLE Frequente  (idPiste NUMBER (10),
                        idSkieur NUMBER (10),
                        PRIMARY KEY (idSkieur, idPiste),
                        FOREIGN KEY (idSkieur) REFERENCES Skieur
                             ON DELETE CASCADE,
                        FOREIGN KEY (idPiste) REFERENCES Piste
                             ON DELETE CASCADE);
\end{verbatim}

La table {\tt Frequente} est un exemple standard
d'association n-m.

\smallskip
{\bf D\'ependances fonctionnelles (3 pts)}


\begin{enumerate}
   \item Oui (un client est dans un seul h\^otel).

    \item Bien s\^ur, puisque rien ne l'interdit !

    \item Idem.

\end{enumerate}


\smallskip
{\bf Requ\^etes ({\bf 6 pts})}

Voici les requ\^etes SQL.

\begin{verbatim}
SELECT nomPiste 
FROM piste 
WHERE couleur = 'Rouge'

SELECT DISTINCT nomHotel 
FROM  Chambre
WHERE capacite >= 6

SELECT nomSkieur, nomPiste
FROM   Skieur s, Frequente f, Piste p
WHERE  s.id = p.idSkieur
AND    p.id = f.idPiste

SELECT nomSkieur
FROM   Skieur s, Hotel h
WHERE  s.nomHotel = h.nomHotel
AND    nomGerant = 'Logeais'

SELECT nomSkieur
FROM   Skieur
WHERE  NOT EXISTS (SELECT 'x' FROM Frequente
                   WHERE   idSkieur = id)

SELECT nomPiste
FROM   Piste p, Frequente f, Skieur s, Hotel h
WHERE  p.couleur = 'Noire'
AND    p.id = f.idPiste
AND    f.idSkieur = s.id
AND    h.nomHotel = s.nomHotel
AND    nomGerant = 'Toirac'
\end{verbatim}

Les requ\^etes en alg\`ebre se d\'eduisent ais\'ement.

\smallskip
{\bf Requ\^etes SQL sur le Zoo ({\bf 4 pts})}
 
\begin{verbatim}
SELECT  nom
FROM    Animal
WHERE   espece IS NOT NULL;

SELECT COUNT(*)
FROM Animal
WHERE   emplacement = 35;

SELECT a1.nom, a2.nom
FROM   Animal a1, Animal a2
WHERE  a1.emplacement = a2.emplacement
AND    a1.nom < a2.nom;

SELECT emplacement, COUNT(*)
FROM   Animal
GROUP BY emplacement
HAVING COUNT(*) >= 4;

\end{verbatim}
\end{document}